\section{The Indirect Branch Problem}

As detailed in Section \ref{sub:shm}, AFL++ relies on a highly optimised, fixed-size shared memory bitmap (typically 64KB, although resizeable) to track execution paths. 

Indirect branches --- such as function pointers, virtual method dispatches in C++, and complex switch-case statements --- introduce a dynamic $1$-to-$N$ mapping into the CFG. A single indirect call site $A$ does not resolve to a single destination $B$. Instead, depending on the program state, $A$ can transition to one of many targets during execution.

If these indirect branches are instrumented using the standard edge-tracking mechanism, every new target $B_i$ generates a new edge $A \rightarrow B_i$.

\paragraph{The Collision Problem} Mapping these high-volume dynamic edges into the same fixed-size map used for direct edges can significantly increase the overall density of the coverage map. A denser map naturally increases the probability of collisions. 

When a collision occurs, it causes "coverage blindness". The fuzzer might observe an execution that reaches a novel path, but if that edge hashes to a byte already saturated by an indirect branch, the fuzzer will incorrectly assume the path is redundant. Consequently, viable inputs are discarded.

\paragraph{Collision-Free Instrumentation Issues} This problem persists even when utilizing collision-free instrumentation modes, such as LLVM PCGUARD (Section \ref{sub:llvmmode}). These modes achieve zero collisions by assigning a unique ID to every basic block and edge discovered at compile time, sizing the coverage map exactly to the number of known edges.

Because indirect branches resolve at runtime, they circumvent the static counting. If unpredictable dynamic edges were naively introduced into a collision-free map, they would risk breaking the zero-collision guarantee.

Simply extending the coverage map to accommodate all dynamic edges would risk allocating an excessively large map, introducing performance bottlenecks and causing CPU cache thrashing.

\subsection{Requirements for a Solution}

To accurately track indirect branches without polluting the primary coverage map or degrading execution speed, an alternative data structure is required. This structure must satisfy two conflicting constraints: 
\begin{itemize}
    \item \textit{High capacity for dynamic edges:} It must be able to absorb a high volume of unpredictable, runtime-generated edges.

    \item \textit{Strict memory and performance bounds:} It must maintain a small memory footprint to avoid cache thrashing and preserve fuzzing throughput.
\end{itemize}

These requirements point towards a probabilistic data structure capable of tracking set membership with a fixed size and a predictable collision rate --- a role perfectly suited for Bloom filters.